
\chapter*{Quarta di copertina}
Questo è il primo di tre volumi dedicati alla meccanica della trave. Vi si sviluppano la cinematica e la statica del corpo rigido e dei sistemi di corpi rigidi, specializzate al caso della trave e dei suoi sistemi --- travature a collegamento eterogeneo, sistemi reticolari, telai piani.

L'impostazione muove da una scelta precisa: il modello rigido non è una semplificazione propedeutica al deformabile, ma il fondamento cinematico e statico-vincolare comune a tutti i modelli strutturali. Studiarlo nella sua generalità --- attraverso la dualità cinematica-statica e il principio dei lavori virtuali, che ne sono il filo conduttore --- significa acquisire la grammatica con cui si costruisce qualunque teoria strutturale successiva. I limiti del modello rigido --- le indeterminazioni statiche, le impossibilità cinematiche --- sono studiati come tali: sono il punto da cui muoverà il volume successivo, dedicato alla trave deformabile.

Il volume si rivolge agli studenti dei corsi di laurea in Ingegneria civile (corso di Scienza delle Costruzioni), in Ingegneria edile-architettura (corso di Statica) e in Architettura (corso di Meccanica delle Strutture), e dei corsi affini, oltre che al lettore in autonomia che voglia rifondare in modo organico la propria comprensione del modello rigido. I prerequisiti sono limitati ai corsi del primo anno: calcolo, geometria analitica, nozioni elementari di algebra lineare, e una familiarità con la meccanica del corpo rigido nei suoi termini fisici.